\begin{chapintro}
統合した値が同じ 0.73 でも、試験の結果がそろっているときと、大きく食い違っているときでは、意味が違います。試験間のばらつき（異質性）を測る三つの数と、予測区間の読み方を説明します。続いて、ばらつきの原因を特徴ごとに探すサブグループ解析を降圧薬 A の例で実行します。原因になっている試験を探す Baujat プロットと、連続した特徴で探すメタ回帰は発展の節に置いてあり、初読では飛ばせます。
\end{chapintro}

\section{ばらつきを測る三つの数}
\label{sec:06-1}

\kkref{sec:05-6} の出力には、統合した効果の下に Q、I²、τ² という三つの数が並んでいました。どれも試験間のばらつきに関係する数ですが、測っているものが違います。

\begin{description}
\item[Q 統計量] 各試験の推定値が統合した値からどれだけ離れているかを、重みを付けて 2 乗して足した値です（重みは \kkref{sec:05-1-3} で説明した分散の逆数です）。各試験の対数リスク比を $y_i$、その重みを $w_i$、統合した値を $\bar{y}$ とすると、$Q = \sum_i w_i (y_i - \bar{y})^2$ と書けます（$\sum_i$ は全試験を足すという記号です）。試験の結果のばらつきが偶然だけなら、Q はおおよそ自由度（試験数 − 1）くらいの値になります。自由度は t 検定で使うものと同じ考え方で、この場合は試験数から 1 を引いた数です。降圧薬 A では Q = 18.30 で、自由度 9 のおよそ 2 倍でした。Q を使った異質性の検定の p 値は 0.032 です。
\item[I²] 試験の結果の散らばりのうち、偶然のぶれでは説明できない部分の割合です。Q と自由度から（Q − 自由度）÷ Q で計算し（Higgins 2002）、降圧薬 A では（18.30 − 9）÷ 18.30 で 51\% になります。
\item[τ²] 真の効果の散らばりそのものの大きさです。ランダム効果モデルの研究間分散（\kkref{sec:05-4}）で、効果指標と同じ目盛（ここでは対数リスク比）で表されます。降圧薬 A では τ² = 0.034、その平方根の τ は 0.18 でした。
\end{description}

異質性の検定は、試験数が少ないと差を検出できません。10 試験程度では、本当にばらつきがあっても p 値が 0.05 を下回らないことがよくあります。そのため、p 値が 0.05 以上でも「異質性はない」とは結論できません。有意水準を 0.10 にする慣習もありますが、検定の結果だけで判断しないほうが安全です。

I² は割合なので、効果の散らばりの大きさそのものは表しません。

\begin{hattenbox}
同じ τ² でも、一つ一つの試験が大きく精度が高いほど偶然のぶれが小さくなるので、I² は大きくなります。たとえば第 2 部に出てくる検査の研究のように、人数の多い研究を集めると、臨床的には小さな違いでも I² が 90\% を超えることがあります。
\end{hattenbox}

I² を 25\%、50\%、75\% で「小さい、中くらい、大きい」と区切る目安が使われることがありますが（Higgins 2003）、この区切りは臨床的な重大さを表すものではありません。I² には信頼区間も付きます。降圧薬 A の I² は 51\% でしたが、95\% 信頼区間は 0〜76\% と広く、10 試験では異質性の程度をほとんど決められないことを示しています。

\begin{hattenbox}
Cochrane ハンドブック（10.10.2）は、0〜40\%、30〜60\%、50〜90\%、75〜100\% という重なりのある範囲を、それぞれ重要でない可能性、中等度、相当、非常に大きいの大まかな目安として挙げ、閾値で機械的に異質性を判断しないよう注意しています。
\end{hattenbox}

τ² は対数の目盛なので、そのままでは読みにくい数です。τ = 0.18 は、真のリスク比が平均の 0.73 から、おおよそ 0.51〜1.04 の範囲に散らばっていることを表します。この幅を、統合した値の不確かさも含めて表したものが、次に説明する予測区間です。

\section{\hatten ばらつきの原因になっている試験を探す}
\label{sec:06-2}

Q と I² は異質性の大きさを一つの数で示すだけで、どの試験がそれを作っているかは教えてくれません。Baujat（ボジャ）プロットは、試験ごとに 2 つの数を計算して 1 点として置く図です。

\begin{rcode}
baujat(m)
\end{rcode}

\kkfig[.8\linewidth]{x1_baujat.pdf}{図 6.1　降圧薬 A の 10 試験の Baujat プロット。点の横の文字は試験名}

横軸は、その試験が Q にどれだけ寄与しているかです。統合した値から大きく離れていて、しかも精度が高い試験ほど右に来ます。縦軸は、その試験を除いて統合し直したときに統合結果がどれだけ動くかです。右にあるほど異質性の原因になっており、上にあるほど結果への影響が大きいことになります。

図 6.1 では、Dubois 2014 が右上に離れて位置しています。2,500 人ずつの大きな試験で、点推定が 0.92 と統合値の 0.73 から離れているため、Q への寄与も、除いたときの統合結果の動きも大きい試験です。Jensen 2022 と Chen 2012 は Dubois 2014 よりさらに右にあり、Q への寄与はいっそう大きいものの、縦の位置は低く、除いても統合結果はそれほど動きません。Jensen 2022 は統合値より効果が大きい側に、Chen 2012 は逆向き（点推定 1.13）に外れた試験で、外れる向きは反対です。Baujat プロットの既定の計算は固定効果モデルに基づくので、\kkref{sec:07-6} のランダム効果モデルで 1 試験ずつ除く解析の数値とは一致しません。

Baujat プロットは、異質性の原因を探す探索用の図です。右上にある試験を見つけたら、その試験の対象、介入、追跡期間、バイアスのリスクに他の試験と違う点がないかを調べる手がかりにします。外れているという理由だけで、その試験を解析から除く根拠にはなりません。除くかどうかは、計画書に書いた基準で決めます。

\section{予測区間と信頼区間}
\label{sec:06-3}

信頼区間の意味が分かりにくいときは \kkref{sec:01-1} を、τ² が分かりにくいときは \kkref{sec:05-4} を読み直してください。

ランダム効果モデルの 95\% 信頼区間 0.61〜0.86 は、真の効果の平均がどこにあるかを表す区間です。試験を増やしていけば、この区間はいくらでも狭くなります。一方で、試験ごとの真の効果は τ だけ散らばっているので、平均が正確に分かっても、ある一つの場面での効果は平均からずれます。

予測区間は、新しい試験を一つ行ったとき（あるいは自分の施設で使ったとき）に、その場面での真の効果がどこに来そうかを表す区間です（Riley 2011）。平均の不確かさと、τ² で表される真の効果の散らばりの両方を含むので、信頼区間より必ず広くなります。降圧薬 A の予測区間は 0.46〜1.15 でした。予測区間は meta の出力にすでに含まれています（\kkref{sec:05-6} で prediction = TRUE としました）。二つの区間の違いを表 6.1 に示します。

\begin{kkcenter}
\kkfit{\begin{tabular}{lll}
\toprule
 & 信頼区間 & 予測区間 \\
\midrule
何の区間か & 真の効果の平均がある範囲 & 新しい一つの試験の真の効果が来そうな範囲 \\
試験を増やすと & 狭くなる & τ の分だけ広いまま残る \\
降圧薬 A & 0.61〜0.86 & 0.46〜1.15 \\
使う場面 & 平均の効果があるかを見る & 自分の場面で効果が期待できるかを見る \\
\bottomrule
\end{tabular}}
\kkcap{表 6.1　信頼区間と予測区間の違い（降圧薬 A のランダム効果モデル）}
\end{kkcenter}

予測区間が 1 をまたいでいるので、降圧薬 A の平均の効果は確かでも、試験を行う場面によっては効果がほとんどない、あるいは逆向きの場面もありうると読みます。メタ解析の結果を診療ガイドラインや施設の方針に使うときには、信頼区間だけでなく予測区間も報告することが勧められています（IntHout 2016）。

\subsection{\hatten 予測区間の計算と作図}
\label{sec:06-3-1}

図 6.2 は、信頼区間と予測区間の違いを、試験ごとの真の効果の分布と並べて描いたものです。

\begin{rcode}
mu  <- m$TE.random        # 真の効果の平均（対数リスク比）
tau <- sqrt(m$tau2)       # 真の効果のばらつき
x <- seq(log(0.35), log(1.6), length.out = 400)
plot(exp(x), dnorm(x, mu, tau), type = "l", log = "x",
     ylim = c(-0.9, 2.3), yaxt = "n", ylab = "",
     xlab = "リスク比（対数目盛）")
abline(v = 1, lty = 2)
segments(exp(m$lower.random), -0.3,    # 信頼区間
         exp(m$upper.random), -0.3, lwd = 5)
segments(exp(m$lower.predict), -0.7,   # 予測区間
         exp(m$upper.predict), -0.7, lwd = 5)
\end{rcode}

\kkfig[.85\linewidth]{x1_pi_ci.pdf}{図 6.2　試験ごとの真のリスク比の分布（平均 0.73、τ = 0.18）と、95\% 信頼区間（0.61〜0.86）、95\% 予測区間（0.46〜1.15）}

釣り鐘の曲線は、試験ごとの真のリスク比が平均 0.73、τ = 0.18（対数目盛で）の分布に従うという、ランダム効果モデルの考え方を描いたものです。この分布の中央 95\% は、平均の対数 ± 1.96τ を比に戻した 0.51〜1.04 にあたります。横軸はリスク比の対数目盛で、点線が 1（差なし）です。上の太い線が 95\% 信頼区間で、この分布の平均がどこにあるかの不確かさを表し、1 をまたぎません。下の太い線が 95\% 予測区間で、これから行う一つの試験の真の効果がどこに来そうかを表します。分布そのものの広がりを含むので、信頼区間よりずっと長く、右端が 1 を越えます。

図の曲線は平均と τ をそのまま使った近似で、予測区間の端とは一致しません。予測区間は平均の推定の不確かさも含めて t 分布で計算しており、10 試験では τ そのものの推定も不安定なので、結果は目安として扱います。

\section{サブグループ解析}
\label{sec:06-4}

ばらつきがあると分かったら、次はその原因を探します。最も単純な方法がサブグループ解析で、試験をある特徴で群に分け、群ごとに統合して比べます。降圧薬 A では、計画の段階で「高齢者を対象にした試験とそれ以外で効果が違うか」を調べると決めていたとします。

\begin{rcode}
m_sub <- update(m, subgroup = pop)
m_sub
\end{rcode}

\begin{rout}
Results for subgroups (random effects model):
              k     RR           95% CI   tau^2
pop = 高齢者  5 0.6470 (0.4606 to 0.9089) 0.0856
pop = 一般    5 0.8221 (0.7139 to 0.9467) <0.0001

Test for subgroup differences (random effects model):
                  Q d.f. p-value
Between groups 1.63    1  0.2020
\end{rout}

高齢者を対象にした 5 試験ではリスク比 0.65（0.46〜0.91）、一般の成人の 5 試験では 0.82（0.71〜0.95）でした。点推定だけ見ると高齢者で効果が大きいようですが、二つの群の差の検定は p = 0.20 で、差があるとは言えません。図 6.3 では、高齢者の群の中でも Chen 2012 だけが 1 より右にあり、群の中のばらつき（I² 60\%）が大きいことが分かります。

\kkfig[.92\linewidth]{p1_forest_sub.pdf}{図 6.3　対象集団で分けたフォレストプロット（架空のデータ）}

\begin{hattenbox}
同じ比較を固定効果モデルで行うと、差の検定は Q = 4.01、p = 0.045 となり、有意になります。固定効果モデルは群の中の真の効果のばらつきを無視するので、群の平均の不確かさを小さく見積もり、群の差を検出しやすくなるからです。Cochrane ハンドブック（10.11.3.1）も、サブグループ間の差の検定はランダム効果モデルに基づくもののほうが望ましいとしています。どちらのモデルで検定するかによって結論が変わるこの例は、サブグループ解析の結果がもろいことを表しています。
\end{hattenbox}

サブグループ解析を読むときは、次の点を確かめます。

\begin{description}
\item[事前に決めていたか] 結果を見てから分け方を探すと、偶然の差がいくらでも見つかります。計画書に書いたサブグループかどうかを確かめます。
\item[差の検定を見ているか] 一方の群だけ有意で他方が有意でないことは、群の間に差があることを意味しません。群の差そのものの検定（Test for subgroup differences）を見ます。
\item[試験の中の比較か、試験の間の比較か] 高齢者の試験と一般の試験は、対象の年齢以外にも、実施した国、時期、併用治療などが違うかもしれません。サブグループ解析で比べているのは試験どうしであり、ランダム化で保証された比較ではないので、観察研究と同じく交絡を受けます。年齢による効果の違いを本当に知りたいなら、各試験の中で年齢別の結果を集めるほうが信頼できます。
\item[群の試験数] 1 群に数本しか試験がないと、群ごとの推定も差の検定も不安定です。
\end{description}

バイアスのリスクで分けたサブグループ（低 6 試験 0.74、懸念あり 3 試験 0.63、高 1 試験 0.50。差の検定 p = 0.67）も同じように計算できます。こちらは \kkref{sec:07-6} の感度分析で扱います。

発展の \kkref{sec:06-2} で紹介した Baujat プロットのように、外れた試験を見つけて異質性の原因の手がかりにすることはできますが、外れているというだけでその試験を解析から除く根拠にはなりません。除くかどうかは、計画書に書いた基準で決めます。

対象の平均年齢や投与量のような連続した特徴で効果の違いを調べる方法をメタ回帰と呼び（発展の \kkref{sec:06-5}）、\kkref{sec:14-2} では検査のカットオフを共変量にして同じ考え方を使います。サブグループ解析もメタ回帰も、説明に使う特徴 1 つあたり少なくとも 10 試験が必要とされています（介入の Cochrane ハンドブック）。降圧薬 A のように試験が 10 本前後しかないときは、事前に決めた 1 つの特徴にとどめ、結果は仮説の提示として扱います。また、これらが扱うのは試験の平均の特徴なので、試験単位の関連を個人の関連と読み替える誤り（生態学的誤謬）に注意が要ります。たとえば対象の平均年齢と効果に関連があっても、各試験の中で高齢の人ほど効果が大きいとは限りません。

\section{\hatten メタ回帰}
\label{sec:06-5}

サブグループ解析を、連続した特徴（対象の平均年齢、試験の年、投与量など）に広げたものがメタ回帰です。試験ごとの効果の大きさを、試験の特徴で説明する回帰分析で、点の大きさが重みに対応します。回帰直線の傾きは、特徴が 1 単位増えたときの効果の変化を表し、切片は特徴が 0 のときの効果を表します。

降圧薬 A の 10 試験では、メタ回帰を行う目安（サブグループ解析の節の終わりに述べました）にぎりぎり届く程度です。そこで、13 試験ある実データを使います。R の metadat パッケージに入っている dat.bcg は、結核を予防する BCG ワクチンの 13 の RCT をまとめたデータです（Colditz 1994）。メタ解析の教科書で繰り返し使われてきた例で、試験を行った土地の緯度（ablat）が記録されています。

\begin{rcode}
library(metafor)   # meta と並ぶメタ解析のパッケージ
dat <- escalc(measure = "RR",
              ai = tpos, bi = tneg,   # ワクチン群
              ci = cpos, di = cneg,   # 対照群
              data = dat.bcg)
r0 <- rma(yi, vi, data = dat, method = "REML")
r1 <- rma(yi, vi, mods = ~ ablat, data = dat,
          method = "REML")
r1
\end{rcode}

escalc は試験ごとの対数リスク比（yi）とその分散（vi）を計算する関数で、rma がランダム効果モデルの当てはめです。mods に説明に使う特徴を書くとメタ回帰になります。まず緯度を入れないモデル（r0）の結果は、リスク比 0.49（0.34〜0.70）、I² = 92\%、予測区間 0.15〜1.55 で、効果は試験によって大きく散らばっていました。

これを試験ごとに見るには、フォレストプロットを緯度の順に並べます。metagen は、試験ごとの効果とその標準誤差（分散の平方根）をそのまま統合する関数です。

\begin{rcode}
mb <- metagen(yi, sqrt(vi), data = dat, sm = "RR",
              studlab = paste(author, year),
              common = FALSE)
forest(mb, sortvar = ablat,           # 緯度の順に並べる
       leftcols = c("studlab", "ablat"))
\end{rcode}

\kkfig[.9\linewidth]{x1_bcg_forest.pdf}{図 6.4　BCG ワクチンの 13 試験のフォレストプロット。緯度（左から 2 列目）の低い順に並べた}

図 6.4 の横軸は対数目盛のリスク比で、1 より左ならワクチン群で結核が少ないことを表します。最下段の菱形がランダム効果モデルで統合した結果で、リスク比 0.49（0.34〜0.70）です。しかし I² は 92\% と非常に大きく、統合した値は個々の試験の結果をあまり代表していません。試験の行を上から順に追うと、緯度が低い（上の）試験では区間が 1 を含むか 1 に近く、緯度が高い（下の）試験ほど区間が 1 から離れて左に寄る傾向が見えます。試験の並べ方を変えるだけで、ばらつきに並び方（パターン）があることが分かります。この並び方を数にするのが、次のメタ回帰です。

緯度を入れたモデル（r1）の出力の要点を抜き出します。

\begin{rout}
tau^2 (estimated amount of residual heterogeneity): 0.0764
I^2 (residual heterogeneity / unaccounted variability): 68.39%
R^2 (amount of heterogeneity accounted for):            75.62%

Test of Moderators (coefficient 2):
QM(df = 1) = 16.3571, p-val < .0001

         estimate      se     zval    pval    ci.lb    ci.ub
intrcpt    0.2515  0.2491   1.0095  0.3127  -0.2368   0.7397
ablat     -0.0291  0.0072  -4.0444  <.0001  -0.0432  -0.0150
\end{rout}

ablat の係数（傾き）−0.029 は、緯度が 1 度上がるごとに対数リスク比が 0.029 下がる、つまりワクチンの効果が大きくなることを表します。切片 0.25 は緯度 0 度のときの対数リスク比です。R² の 76\% は、試験間の真の効果のばらつき（τ²）のうち 76\% を緯度で説明できたという意味で、τ² は 0.31 から 0.076 に減りました。それでも残りのばらつき（残差の I² 68\%）は小さくありません。緯度 20 度、40 度、55 度での予測リスク比は、それぞれ 0.72（0.55〜0.93）、0.40（0.32〜0.51）、0.26（0.17〜0.38）です。

試験数が少ないメタ回帰では、発展の \kkref{sec:05-7-1} で説明した Hartung-Knapp 調整を係数の検定と区間に使うことが多く、metafor では rma() に test = "knha" を足すだけです。

\begin{rcode}
r1k <- rma(yi, vi, mods = ~ ablat, data = dat,
           method = "REML", test = "knha")
\end{rcode}

調整すると、ablat の係数は同じ −0.029 のまま、p 値は 0.0046、95\% 信頼区間は −0.047〜−0.011 となり、区間は広がりますが、緯度が高いほど効果が大きいという結論は変わりません。

\kkfig[.82\linewidth]{p1_bcg_bubble.pdf}{図 6.5　BCG ワクチンの 13 試験のメタ回帰。円の大きさは重み、帯は回帰直線の 95\% 信頼区間}

赤道に近い地域では、環境中の抗酸菌にもともと触れているためにワクチンの上乗せ効果が小さくなる、という説明が考えられてきました。ただし、メタ回帰で見つかった関連は、試験の間の観察にすぎません。緯度の違う試験は、実施した年代、割り付けの方法、対象の年齢なども違っており、緯度がそれらの代わりに関連を示している可能性があります。

\section{この章のまとめ}
\label{sec:06-6}

異質性は三つの数で見ます。Q と検定は試験数が少ないと差を検出できず、I² は偶然では説明できない散らばりの割合であって散らばりの大きさではなく、τ² が真の効果の散らばりそのものです。降圧薬 A の予測区間 0.46〜1.15 は 1 をまたいでおり、平均の効果が明らかでも、場面によって効果が小さくなりうることを示しています。サブグループ解析は試験どうしの比較なので交絡を受け、事前に決めた少数の特徴にとどめます。降圧薬 A の対象集団による違いは、ランダム効果モデルでは有意ではありませんでした。発展の節で扱った Baujat プロットとメタ回帰は、異質性の原因を探す手がかりであり、除外や因果の根拠にはなりません。次の章では、集めた試験そのものに偏りがないか、結果がどこまで確かかを調べます。
